\documentclass[11pt]{article}

\usepackage{amsmath, amssymb}
\usepackage{amsthm}

\newtheorem{lemma}{Lemma}

\title{Appendix for ``Estimating Representative Causal Effects with Double Machine Learning''}
\author{Apoorva Lal \and Winston Chou}
\date{}

\begin{document}

\maketitle

\appendix

\section{Detailed Derivations}

\subsection{Derivation of the Conditional Causal Derivative}
\label{sec:poisson-derivation}
Note that, in our stylized example, $T_i$ is $\text{Poisson}(\lambda_j)$ conditional on $X_i = j$.  Therefore, its probability mass function is given by:
\begin{equation}
    P(T_i = t | X_i = j) = e^{-\lambda_j} \frac{\lambda_j^t}{t!}, t = 0, 1, 2, \ldots,
\end{equation}
and the conditional expectation of the causal derivative $\frac{1}{T_i + 1}$ given $X_i = j$ is:
\begin{equation}
    \sum_{t = 0}^\infty \frac{1}{t + 1} e^{-\lambda_j} \frac{\lambda_j^t}{t!} = \sum_{t = 0}^\infty e^{-\lambda_j} \frac{\lambda_j^t}{(t + 1)!}.
\end{equation}

Define $s := t + 1$.  Then we can rewrite the above as:
\begin{eqnarray}
\sum_{s = 1}^\infty e^{-\lambda_j} \frac{\lambda_j^{s - 1}}{s!} &=& \frac{e^{-\lambda_j}}{\lambda_j} \sum_{s = 1}^\infty \frac{\lambda_j^s}{s!} \\
&=& \frac{(e^{\lambda_j} - 1)}{\lambda_j e^{\lambda_j}} \\
&=& \frac{1 - e^{-\lambda_j}}{\lambda_j}.
\end{eqnarray}
where the second-to-last equality uses the identity $\sum_{s = 0}^\infty \frac{x^s}{s!} = e^x$.


\subsection{Derivation of $T_i^*$}

\label{sec:poisson-tstar-derivation}

Within a given stratum $X_i = j$, $h(X_i) = E[T_i | X_i] = \lambda_j$ by assumption.  We want to derive the point $T_i^*$ between $T_i$ and $\lambda_j$ at which we are evaluating the derivative $f'(t) = 1 / (t + 1)$.  We will condition on $X_i$ throughout.

By the mean value theorem, there is some $T_i^*$ between $T_i$ and $\lambda_j$ for which:
\begin{equation}
    f(T_i) - f(\lambda_j) = f'(T_i^*) (T_i - \lambda_j).
\end{equation}

Plugging in the definition of $f$ and its derivative, we write this as:
\begin{equation}
    \log(T_i + 1) - \log(\lambda_j + 1) = \frac{1}{T_i^* + 1} (T_i - \lambda_j).
\end{equation}
Solving for $T_i^*$ yields:
\begin{equation}
    T_i^* = \frac{T_i - \lambda_j}{\log \frac{T_i + 1}{\lambda_j + 1}} - 1.
\end{equation}
Note that, when $T_i = \lambda_j$, this quantity is undefined (in which case we just set $T_i^* = T_i$).

In our simulations, we estimate the theoretical plim of RORR by taking many draws of $T_i$, plugging them into this formula, and using the resulting stratum means to estimate the RORR plim.

\section{Choosing the Number of Segments}

\label{app:guidance}

Choosing the number of segments $K$ into which to coarsen the treatment involves
the usual considerations of bias and variance. Larger values of $K$ reduce bias, but they also reduce the sample size in each segment, increasing variance.  Increasing $K$ can also make results less interpretable.  In our day-to-day work, we find that a relatively small number of bins---for example, dividing users into low, medium, and high usage segments---is sufficient to uncover meaningful heterogeneity in the dose-response function. 

Here, we give a highly informal justification for choosing a relatively small number of bins, on the order of $N^{1/5}$.  In practice, researchers should assess the sensitivity of their results to different choices of $K$.

We assume the usual conditions for the coarsened AIPW estimator, i.e., that the treatment is conditionally unconfounded and is continuously distributed on a compact interval.  As in the main text, we decompose the error of the coarsened AIPW estimator as
\begin{align}
\nonumber 
& \widehat{\psi} - E[f'(T_i)] \\
& \quad =
\underbrace{\left\{
\sum_{k=1}^{K-1}
w_k
\left(
\frac{\psi_{k+1}-\psi_k}{\ell}
\right)
-E[f'(T_i)]
\right\}}_{\text{coarsening bias}} \nonumber
\\
& \qquad +
\underbrace{
\sum_{k=1}^{K-1}
w_k
\left(
\frac{
(\hat\psi_{k+1}-\psi_{k+1})
-
(\hat\psi_k-\psi_k)
}{\ell}
\right)
}_{\text{estimation error}}.
\end{align}

Theorem 1 showed that the coarsening bias is
\begin{equation}
\sum_{k=1}^{K-1}
w_k
\left(
\frac{\psi_{k+1}-\psi_k}{\ell}
\right)
-
E[f'(T_i)]
=
O(\ell)
=
O(K^{-1}),
\end{equation}
and therefore the squared coarsening bias is $O(K^{-2})$.

Next, we consider sampling variation. With cross-fitting and the usual AIPW
regularity conditions, we have that
\begin{equation}
\hat\psi_k-\psi_k
=
\frac{1}{N}\sum_{i=1}^N \phi_k(O_i)
+
R_{k,N},
\end{equation}
where $\phi_k$ is the AIPW influence function and $R_{k,N}$ is
a higher-order remainder due to nuisance estimation.  Under smoothness of $p(t\mid x)$, $\pi_k = O(\ell)$.  Therefore, the variance of the bin-level AIPW-estimated mean is
\begin{equation}
\operatorname{Var}(\hat\psi_k)=O\left(\frac{1}{N\ell}\right)
=
O\left(\frac{K}{N}\right).
\end{equation}

It remains to account for the remainder due to nuisance estimation. Let
\begin{equation}
r_{N,K}=\max_{1\le k\le K}|R_{k,N}|.
\end{equation}
Plugging this into the estimation error expression (and ignoring telescoping for a conservative bound) yields
\begin{equation}
\left|
\sum_{k=1}^{K-1}
w_k
\left(
\frac{R_{k+1,N}-R_{k,N}}{\ell}
\right)
\right|
=
O_p(Kr_{N,K}),
\end{equation}
so its squared contribution is $O_p(K^2r_{N,K}^2)$.

Combining these pieces, the coarsened AIPW estimator has MSE
\begin{equation}
\operatorname{MSE}(\widehat\psi)
=
O(K^{-2})
+
O\left(\frac{K}{N}\right)
+
O_p(K^2r_{N,K}^2).
\end{equation}
A typical assumption for the nuisance remainder is that it is $o_p((N\ell)^{-1/2})$, so a conservative bound sets
\begin{equation}
r_{N,K}
=
O_p\left(\left(\frac{N}{K}\right)^{-1/2}\right)
=
O_p\left(\sqrt{\frac{K}{N}}\right).
\end{equation}
Under this assumption, the nuisance-remainder contribution is
\begin{equation}
O_p(K^2r_{N,K}^2)
=
O_p\left(\frac{K^3}{N}\right),
\end{equation}
and so
\begin{equation}
\operatorname{MSE}(\hat\Psi_K)
=
O(K^{-2})
+
O\left(\frac{K^3}{N}\right),
\end{equation}
up to lower-order terms.  The $K$ that minimizes this is $K^*=O(N^{1/5})$

Again, this derivation is heuristic and should be treated as a starting point. Researchers should assess the sensitivity of their results to several values of $K$ while also scrutinizing the underlying assumptions of continuity and positivity of the treatment.

\section{Lemma~\ref{lemma:midpoint}}

\label{app:midpoint}

\begin{lemma}
    \label{lemma:midpoint}
    Assume $f$ is twice continuously differentiable with $|f''(t)|\le M$ and that $p(t\mid x)$ is twice continuously differentiable with uniformly bounded derivatives. Also assume positivity:
    \[
    \int_{S_k}p(t\mid x)dt \ge \epsilon \ell
    \]
    uniformly in $k$ and $x$.

    Then for every $x$ and uniformly in $k$,
    \begin{equation}
        \int f(t) r_k (t|x) dt = f(\overline{t}_k) + O(\ell^2),
    \end{equation}
    where, as in the main text, $\overline{t}_k = \frac{t_{k+1} + t_k}{2}$, $\ell = t_{k + 1} - t_k$, and $r_k(t|x) = \frac{p(t|x) 1(t \in S_k)}{\int_{S_k} p(t|x)dt}.$
\end{lemma}

\begin{proof}
    Since $f$ is twice continuously differentiable,
    \begin{align}
    \nonumber \int f(t) r_k(t|x) dt &= f(\overline{t}_k) + f'(\overline{t}_k)\int(t - \overline{t}_k) r_k(t|x) dt \\
    & \quad + \frac{1}{2} \int f''(\tilde{t}) (t - \overline{t}_k)^2 r_k(t|x) dt
    \end{align}
    for some $\tilde{t}$ between $t$ and $\overline{t}_k$.

    Since $f''$ is bounded and $(t - \overline{t}_k)^2 \le (\ell/2)^2$,
    \begin{equation}
    \left|\frac{1}{2} \int f''(\tilde{t}) (t - \overline{t}_k)^2 r_k(t|x) dt\right| \le C\ell^2.
    \end{equation}
    Therefore, it remains to show that the other error term $\int(t - \overline{t}_k) r_k(t|x) dt = O(\ell^2)$.

    We can rewrite the numerator of this term as
    \begin{equation}
    \int_{-\ell/2}^{\ell/2} u p(\overline{t}_k + u|x)du,
    \end{equation}
    where $u := t - \overline{t}_k$.  Since $p$ is twice continuously differentiable, this is equal to
    \begin{equation}
    \int_{-\ell/2}^{\ell/2} u \left[p(\overline{t}_k|x) + up'(\overline{t}_k|x) + \frac{1}{2} u^2 p''(\xi|x)\right]du
    \end{equation}
    for some $\xi$ between $t$ and $\overline{t}_k$.  The first term integrates to zero because $\int_{-\ell/2}^{\ell/2} udu = 0$. Therefore, the numerator depends only on second- and higher-order powers of $u$.  Since $u^2 \le \ell^2$, the numerator is $O(\ell^3)$.

    Turning to the denominator and defining $u$ as before, we have
    \begin{equation}
        \int_{S_k} p(t|x) dt = \int_{-\ell/2}^{\ell/2} p(\overline{t}_k + u|x)du.
    \end{equation}
    By the same expansion, this is $p(\overline{t}_k|x) \ell + O(\ell^2)$.  Putting the numerator and denominator together, we have $O(\ell^3) / (p(\overline{t}_k|x) \ell + O(\ell^2)) = O(\ell^2)$ by positivity, which gives the result.
\end{proof}



\end{document}
